\chapter{Conclusions}\label{ch:conclusions}

Ce corpus fournit une documentation située sur les manières de nommer les morts, de dater des événements et de représenter les relations dans les inscriptions chrétiennes de Rome entre 340 et 403. Il ne suffit ni à reconstituer la population chrétienne de la Ville, ni à mesurer l’adhésion doctrinale des commanditaires. Sa valeur tient à la possibilité de relier chaque proposition au texte transmis, à son second témoin et à son statut documentaire.

La confrontation systématique de deux bases épigraphiques indépendantes a transformé l’état du dossier. Elle a fait apparaître des formules de date que le premier témoin omettait (\nnums{84}, 214, 268, 388, 447…), des lignes absentes de l’un ou de l’autre (\nnums{306}, 378), des restitutions divergentes qui commandent l’interprétation (\nnums{73}, 447), des erreurs de saisie (dates de l’EDB aux \nnums{206} et 278 ; table des personnes de la base de travail), et des données matérielles jusque-là ignorées : supports, dimensions, signes, reproductions. Elle a aussi montré ses limites : deux bases qui concordent ne font pas une édition collationnée.

La chronologie a été rapportée à ses fondements. Naître, mourir, être déposé, acquérir une tombe et faire graver un texte ne sont pas une même opération. Une concordance arithmétique contrôle une relation, non l’ensemble d’une lecture ; une discordance appelle l’examen du témoin et n’autorise pas à normaliser silencieusement le texte. L’année de classement de chaque notice est conservée ; sa solidité est qualifiée.

Le vocabulaire religieux appelle la même discipline. Les désignations de Dieu et du Christ sont rares, la Trinité n’est mentionnée qu’une fois, l’Esprit saint au plus une fois et sur restitution, la résurrection jamais de manière assurée. La foi s’exprime surtout par la paix et le repos, par le chrisme, et par la qualification des personnes — \emph{fidelis}, \emph{neophytus}, \emph{ancilla dei}, \emph{virgo}. Quelques poèmes, souvent issus de milieux aisés ou proches du clergé, développent l’espérance d’une vie éternelle, de la lumière et de la société des saints. Le silence doctrinal du genre ne peut être lu comme le silence d’une communauté.

La contribution du volume est ainsi documentaire et méthodologique avant d’être interprétative : distinguer les niveaux de transmission, publier les règles de compte, confronter les témoins, conserver les états antérieurs et limiter chaque inférence à la résolution des preuves disponibles. Le catalogue, les données et les programmes qui produisent chaque tableau sont publiés avec le livre, afin que toute mesure puisse être recalculée et toute décision contestée. L’établissement philologique complet, qui suppose la collation des éditions imprimées et des monuments, reste une étape distincte ; ce volume en fournit l’instrument et en délimite le programme.
